\section{هدف و مسیر فرمان}
هدف این مرحله، ارسال یک یا چند
\lr{action}
از پایتون، اجرای کنترل‌شده در برد و دریافت پاسخ مرتبط است:

\begin{LTR}
\begin{lstlisting}
command_sender.py -> MQTT control -> ESP8266 -> actuators.cpp
command_sender.py <- MQTT response <- ESP8266 <--- result
\end{lstlisting}
\end{LTR}

فرستنده پیش از
\lr{publish}
روی موضوع پاسخ
\lr{subscribe}
می‌کند. در نتیجه پاسخ سریع برد از دست نمی‌رود. هر فرمان یک
\lr{command\_id}
یکتا دارد و فرستنده فقط پاسخی را می‌پذیرد که همین شناسه را داشته باشد.

\section{موضوع‌های کنترل}
برای دستگاه
\lr{esp8266-01}
دو موضوع استفاده می‌شود:

\begin{LTR}
\begin{lstlisting}
iot/devices/esp8266-01/control
iot/devices/esp8266-01/response
\end{lstlisting}
\end{LTR}

برنامهٔ پایتون فرمان را روی
\lr{control}
منتشر می‌کند. برد پس از پردازش، نتیجه را روی
\lr{response}
می‌فرستد. شناسهٔ دستگاه بخشی از موضوع است؛ بنابراین فرمان یک برد به برد دیگر ارسال نمی‌شود.

\section{ساختار فرمان تک‌عملی}
برای روشن‌کردن چراغ داخلی:

\begin{LTR}
\begin{lstlisting}
{
  "schema_version": 1,
  "command_id": "cmd-001",
  "actions": [
    {"target": "led", "value": true}
  ]
}
\end{lstlisting}
\end{LTR}

\lr{actions}
همیشه یک فهرست است؛ حتی وقتی فقط یک عمل وجود دارد. این تصمیم باعث می‌شود برای چند عمل به قرارداد دیگری نیاز نباشد.

\section{ساختار فرمان چندعملی}
نمونهٔ زیر هم‌زمان چراغ را روشن و دو
\lr{setpoint}
را تغییر می‌دهد:

\begin{LTR}
\begin{lstlisting}
{
  "schema_version": 1,
  "command_id": "cmd-002",
  "actions": [
    {"target": "led", "value": true},
    {"target": "temperature_setpoint", "value": (*@\lr{24.5}@*)},
    {"target": "humidity_setpoint", "value": (*@\lr{60.0}@*)}
  ]
}
\end{lstlisting}
\end{LTR}

حداکثر هشت
\lr{action}
در یک فرمان پذیرفته می‌شود. برای تعداد بیشتر، چند فرمان جدا ارسال کنید. این محدودیت اندازهٔ پیام و حافظهٔ برد را قابل پیش‌بینی نگه می‌دارد.

\section{هدف‌های کنترلی فعلی}
\begin{itemize}
  \item
  \lr{led}
  مقدار منطقی می‌گیرد و چراغ داخلی را کنترل می‌کند؛
  \item
  \lr{heater}
  ،
  \lr{cooler}
  و
  \lr{relay}
  مقدار منطقی می‌گیرند، اما خروجی آن‌ها در تنظیم اولیه غیرفعال است؛
  \item
  \lr{temperature\_setpoint}
  عددی بین
  \lr{-40}
  و
  \lr{125}
  می‌گیرد؛
  \item
  \lr{humidity\_setpoint}
  عددی بین صفر و صد می‌گیرد.
\end{itemize}

مقدارهای
\lr{setpoint}
فعلاً در
\lr{RAM}
هستند و با راه‌اندازی دوبارهٔ برد به مقدار اولیه برمی‌گردند. ذخیرهٔ دائمی فقط زمانی اضافه می‌شود که رفتار بازیابی و فرسایش حافظه طراحی شده باشد.

منطق کنترل خودکار آینده می‌تواند مقدارها را بدون وابستگی به
\lr{MQTT}
از خود ماژول بخواند:

\begin{LTR}
\begin{lstlisting}[language=C++]
const float wantedTemperature = getTemperatureSetpoint();
const float wantedHumidity = getHumiditySetpoint();
\end{lstlisting}
\end{LTR}

\section{پاسخ برد}
پاسخ یک فرمان موفق:

\begin{LTR}
\begin{lstlisting}
{
  "schema_version": 1,
  "device_id": "esp8266-01",
  "command_id": "cmd-002",
  "success": true,
  "results": [
    {"target":"led","success":true,"message":"LED turned on"},
    {"target":"temperature_setpoint","success":true,
     "message":"temperature setpoint updated"},
    {"target":"humidity_setpoint","success":true,
     "message":"humidity setpoint updated"}
  ]
}
\end{lstlisting}
\end{LTR}

فیلد اصلی
\lr{success}
فقط زمانی
\lr{true}
است که همهٔ
\lr{action}ها
موفق باشند.

\subsection{نمونهٔ موفقیت جزئی}
در تنظیم اولیه چراغ داخلی فعال ولی
\lr{relay}
غیرفعال است. اگر هر دو در یک فرمان روشن شوند، چراغ اجرا می‌شود و
\lr{relay}
رد می‌شود:

\begin{LTR}
\begin{lstlisting}
{
  "command_id": "cmd-003",
  "success": false,
  "results": [
    {"target":"led","success":true,"message":"LED turned on"},
    {"target":"relay","success":false,
     "message":"relay output is disabled"}
  ]
}
\end{lstlisting}
\end{LTR}

\lr{action}ها
به‌ترتیب اجرا می‌شوند و نتیجهٔ مستقل دارند. بنابراین فرستنده باید همیشه آرایهٔ
\lr{results}
را بخواند و فقط به مقدار کلی اکتفا نکند.

\section{اجرای فرستندهٔ پایتون}
از پوشهٔ
\lr{code/collector}
و با محیط مجازی فعال اجرا کنید.

روشن‌کردن چراغ داخلی:

\begin{LTR}
\begin{lstlisting}
python command_sender.py esp8266-01 --set led=on
\end{lstlisting}
\end{LTR}

خاموش‌کردن چراغ داخلی:

\begin{LTR}
\begin{lstlisting}
python command_sender.py esp8266-01 --set led=off
\end{lstlisting}
\end{LTR}

تغییر یک
\lr{setpoint}:

\begin{LTR}
\begin{lstlisting}
python command_sender.py esp8266-01 \
  --set temperature_setpoint=(*@\lr{24.5}@*)
\end{lstlisting}
\end{LTR}

ارسال چند
\lr{action}
در یک فرمان:

\begin{LTR}
\begin{lstlisting}
python command_sender.py esp8266-01 \
  --set led=on \
  --set temperature_setpoint=(*@\lr{24.5}@*) \
  --set humidity_setpoint=60
\end{lstlisting}
\end{LTR}

برای
\lr{Windows PowerShell}
به‌جای
\lr{backslash}
از
\lr{backtick}
استفاده کنید:

\begin{LTR}
\begin{lstlisting}
python command_sender.py esp8266-01 `
  --set led=on `
  --set temperature_setpoint=(*@\lr{24.5}@*) `
  --set humidity_setpoint=60
\end{lstlisting}
\end{LTR}

مقادیر
\lr{on/true}
به
\lr{true}
و مقادیر
\lr{off/false}
به
\lr{false}
تبدیل می‌شوند. مقدارهای دیگر باید عدد باشند.

\section{استفاده از فرستنده در کد پایتون دیگر}
رابط برنامه‌نویسی مستقیم برای استفادهٔ آینده در وب یا
\lr{automation}:

\begin{LTR}
\begin{lstlisting}[language=Python]
from command_sender import CommandSender, build_command, parse_action

actions = [
    parse_action("led=on"),
    parse_action("temperature_setpoint=(*@\lr{24.5}@*)"),
]
command = build_command(actions)
response = CommandSender().send("esp8266-01", command)

for result in response["results"]:
    print(result["target"], result["success"], result["message"])
\end{lstlisting}
\end{LTR}

وب‌اپلیکیشن آینده باید همین کلاس را صدا بزند و قرارداد جداگانه‌ای نسازد.

\section{فعال‌سازی ایمن خروجی‌های خارجی}
در فایل
\lr{src/actuators.cpp}
خروجی‌های
\lr{heater}،
\lr{cooler}
و
\lr{relay}
عمداً
\lr{false}
هستند:

\begin{LTR}
\begin{lstlisting}[language=C++]
constexpr bool ENABLE_HEATER_OUTPUT = false;
constexpr uint8_t HEATER_PIN = D1;

constexpr bool ENABLE_COOLER_OUTPUT = false;
constexpr uint8_t COOLER_PIN = D2;

constexpr bool ENABLE_RELAY_OUTPUT = false;
constexpr uint8_t RELAY_PIN = D5;
\end{lstlisting}
\end{LTR}

پیش از تغییر هر گزینه به
\lr{true}
این موارد را بررسی کنید:

\begin{itemize}
  \item پین با فرایند راه‌اندازی برد تداخل نداشته باشد؛
  \item
  \lr{GPIO}
  بار قدرت را مستقیم تغذیه نکند و مدار راه‌انداز یا
  \lr{relay}
  مناسب وجود داشته باشد؛
  \item ولتاژ، جریان، ایزولاسیون، فیوز و حالت قطع ارتباط بررسی شده باشد؛
  \item سطح فعال
  \lr{HIGH/LOW}
  با ماژول واقعی تطبیق داشته باشد؛
  \item خروجی هنگام شروع برنامه ابتدا خاموش بماند.
\end{itemize}

کد اجازه نمی‌دهد
\lr{heater}
وقتی
\lr{cooler}
روشن است فعال شود و برعکس. این قفل متقابل نرم‌افزاری جای حفاظت الکتریکی و ترموستات مستقل را نمی‌گیرد.

\section{افزودن هدف کنترلی تازه}
برای افزودن مثلاً
\lr{pump}
:

\begin{enumerate}
  \item پین، سطح فعال و گزینهٔ فعال‌سازی را در
  \lr{actuators.cpp}
  تعریف کنید؛
  \item در
  \lr{initializeActuators}
  خروجی را در حالت خاموش آماده کنید؛
  \item یک شاخهٔ تازه در
  \lr{applyControlAction}
  اضافه و نوع مقدار را بررسی کنید؛
  \item یک آزمون برای تجزیهٔ فرمان پایتون و یک آزمون واقعی خروجی اضافه کنید؛
  \item هدف کنترلی و محدودیت ایمنی آن را در همین فصل ثبت کنید.
\end{enumerate}

منطق هدف تازه نباید داخل تابع بازگشتی مربوط به
\lr{MQTT}
قرار گیرد؛ تابع بازگشتی فقط
\lr{JSON}
را می‌خواند و ماژول
\lr{actuator}
را صدا می‌زند.

\section{تست و خروجی مورد انتظار}
تست‌های پایتون را اجرا کنید:

\begin{LTR}
\begin{lstlisting}
python -m unittest -v
\end{lstlisting}
\end{LTR}

اکنون در مجموع نه آزمون اجرا می‌شود. ساخت نسخهٔ
\lr{0.3.1}
نیز موفق بوده است. پس از بارگذاری، خروجی سریال مورد انتظار شامل این سطرهاست:

\begin{LTR}
\begin{lstlisting}
Subscribed to control topic: iot/devices/esp8266-01/control
Command response published: {"schema_version":1,...}
\end{lstlisting}
\end{LTR}

برای آزمون اولیه فقط چراغ داخلی را کنترل کنید. آزمون
\lr{heater}،
\lr{cooler}
یا
\lr{relay}
بدون سخت‌افزار قدرت مناسب و حفاظت مستقل انجام نشود.

\section{استقلال کنترل از ارسال داده}
در نسخهٔ
\lr{0.3.1}
پردازش فرمان منتظر زمان انتشار سنسورها نمی‌ماند. تابع
\lr{mqttClient.loop()}
در هر دور برنامه پیش از کار سنسورها اجرا می‌شود. همچنین اگر داده‌ها چند بسته باشند، در هر دور فقط یک بسته منتشر می‌شود و سپس کنترل دوباره به مسیر
\lr{MQTT}
برمی‌گردد.

بنابراین افزایش تعداد سنسورها باعث اجرای یک حلقهٔ طولانی برای انتشار همهٔ بسته‌ها نمی‌شود. در پیاده‌سازی سنسور واقعی نیز نباید داخل
\lr{readSensors}
از تأخیر طولانی استفاده شود. سنسوری که تبدیل زمان‌بر دارد باید شروع اندازه‌گیری و خواندن نتیجه را در دو زمان جدا انجام دهد.
